% Online Supplement F-source: full price regression grid + adaptation detail
% Migrated 2026-06-06 from sec_app04_scope_submission.tex and
% sec_app05_adaptive_deployment_submission.tex during the JLEO compression pass.
% Price is SCOPE only — no damages, overcharge, markup, or causal price effect.
\section{Price Scope and Adaptation: Full Grids}
\label{supp:price}

This section holds the full price regression grid behind the scope reading
of Appendix~F (\ref{app:scope_adaptation_submission}) and the strategic-%
adaptation simulation detail behind \ref{app:strategic_adaptation_submission}.
Every number is descriptive: the price evidence establishes economic
non-neutrality of flagged \emph{environments} and is \emph{not} a damages
estimate, an overcharge, a markup, or a treatment effect on price. The
design holds the procurement environment fixed rather than against a
counterfactual price, so no overcharge can be computed.

%--------------------------------------------------------------------
\subsection{Broad and Segment Price Grids}
\label{supp:price_grids}

The compact sign-reversal decomposition (broad weighting to within-cell
ATT) is reported in Appendix~F (\ref{app:overlap_scope_submission}). This
section holds the full broad-by-modality grid and the quartile/trim segment
grid behind it. In the broad specification, frequent-loser presence is
positively associated with the log negotiated unit price, with a headline
range of $+3.6\%\,\text{--}\,+7.7\%$ across estimators (broad coefficient
$\valBetaBroad$, $p=\valBetaBroadP$; the larger association in the
electronic-auction environment is inconsistent with a pure quorum-filler
reading). These are descriptive associations, not causal effects.

The negative ATT reading survives both modalities (Convite
$\valSignBetaConvATT$, Preg\~ao $\valSignBetaPregATT$) and every
tender-value quartile except the highest; the high-value tail (Q4) is the
only systematically positive cell ($\valBetaQfour$, $p=\valBetaQfourP$),
where coordination costs are spread over larger contracts. Restricting the
overlap to items involving firms named as direct CADE defendants yields a
null price estimate ($\valBetaCADEDirect$, $p=\valBetaCADEDirectP$;
non-direct overlap $\valSignBetaCADENonDirectATT$): the price regressions
do not recover the legal object of the adjudicated cases. The broad
positive imprint is selection into structurally higher-priced cells---among
non-treated items, mean log price rises from $\valSelTestQOne$ to
$\valSelTestQFive$ across cell FL-share quintiles, a $\Delta=\valSelTestDelta$
log-point gap, with a partialled-out coefficient $\valSelTestCoefFullFE$
once all five cell dimensions are absorbed---and the within-cell ATT removes
exactly this component.

%--------------------------------------------------------------------
\subsection{Strategic-Adaptation Simulation Detail}
\label{supp:price_adaptation}

Holding the empirical environment fixed, the participation score is
re-evaluated under simple evasive responses; the exercise is
retrospective, the labels remain adjudication-anchored, and the operating
points are conditional on the validated incumbent ranking. Rotation of
losing-role firms and occasional small wins leave discrimination
essentially unchanged ($\valAUCAdaptRotation$ and $\valAUCAdaptOccWins$
versus baseline $\valAUCBase$), absorbed by the persistence signal; CNPJ
splitting lowers AUC to $\valAUCAdaptCNPJsplit$ by breaking participation
history; threshold-aware capping is the most damaging
($\valAUCAdaptThreshCap$, and $\valAUCAdaptCombined$ combined with
rotation). The hierarchy maps to the cost--recall frontier: a published
threshold invites bunching (consistent with the local density ratio
$\approx 1.06$ at $T=14$ in the profile appendix, where the threshold was
computed ex post and never disclosed) and is the most gameable rule, so
continuous ranks or capacity-constrained top-$k$ queues should be preferred
over a static published cutoff. This is a diagnostic, not a full
equilibrium model of evasion: it measures how fast a static published
cutoff loses its triage value once firms respond to it, and makes no claim
about the durability of any underlying behavioral pattern.
